%% METIS_DRDS.tex
%%
%% METIS data reduction pipeline design document
%%
%% 2020-05-27: initialised from specification document (OC)
%%
\newcommand{\doctitle}{%
  \parbox[t]{\textwidth}{Data Reduction Library\\Design}}
\newcommand{\docnumber}	{E-xxx-AST-MET-xxxx}     % Document number
\newcommand{\issuenumber}{0-1}                % Version
\newcommand{\issuedate}{2020-05-27}              % Date
\newcommand{\workpackage}{Science Pipeline}      % Workpackage

\newcommand{\authorname}{%                       % Authors
  O.~Czoske,\\
  W.~Kausch,\\
  N.~Sabha,\\
  W.~Zeilinger}
\newcommand{\authorsigndate}  {} % date of signature of author

\newcommand{\reviewername}     {%
}       % checked by name
\newcommand{\reviewersigndate} {} % date of signature of checker

\newcommand{\approvername}    {%
}       % name of the approver
\newcommand{\approvalsigndate}{} % date of signature of approver

\newcommand{\releasername}{%
}
\newcommand{\releasesigndate}{}  % date of release

% Title displayed in the header
\newcommand{\shorttitle}{Data Reduction Library Design}

\newif\ifshowSVN \showSVNfalse  % flag to show SVN info instead of Issue info



%%%%%%%%%%%%%%%%%%%%%%%%%%%%%%%%%%%%%%%%%%%%%%%%%%%%%%%%%%%%%%%%%%%%%%%%%%%%%
\documentclass[11pt,oneside,a4paper]{article}

%% Layout  --- TeX commands
\raggedbottom
\topmargin=-9mm
\headsep=0.5cm
\textheight=246mm
\textwidth=164mm
\hoffset=0cm
\oddsidemargin=-2mm
\parindent=0mm
\parskip=0.3em

\renewcommand{\floatpagefraction}{0.98} % separate float page does not work
\renewcommand{\bottomfraction}{1.0}

%% Fonts
\usepackage{mathptmx}
\usepackage[scaled=0.92]{helvet}

\usepackage[pdftex]{graphicx}
\graphicspath{{./figures/}}

\usepackage{fancyhdr}
\usepackage{titlesec}
\usepackage{booktabs}
\usepackage{tabularx}
\usepackage{array}
\usepackage{longtable}
\usepackage{numprint}
\usepackage{placeins}
\usepackage{enumitem}

\usepackage{xcolor}
\usepackage{colortbl}
\definecolor{listingbg}{gray}{0.95}
\definecolor{darkgreen}{rgb}{0.0, 0.7, 0.0}
\definecolor{darkblue} {rgb}{0.0, 0.0, 0.7}
\definecolor{darkred}  {rgb}{0.7, 0.0, 0.0}
\definecolor{darkorange}{rgb}{1.0, 0.49, 0.0}

\usepackage[many]{tcolorbox}
\usepackage{multirow}
\usepackage[autokw=true]{svn-multi}

% does not work with my version of biblatex:
% \usepackage[defernumbers=true,backend=biber,hyperref=true,url=false]{biblatex}
\usepackage[%
   defernumbers=true,
   backend=biber,
   hyperref=true,
   url=false,
   sorting=none]{biblatex}

\usepackage[%
% pagebackref=true,   % does not work with biblatex
   pdfpagelabels=true,
   plainpages=false,
   colorlinks=false]{hyperref}

% does not work with my version:
%\usepackage[xindy,nopostdot,nogroupskip]{glossaries}
\usepackage[xindy]{glossaries}
%\makeglossaries

\usepackage{lastpage}
\usepackage{multirow}
\usepackage[%
   font=small,
   format=plain,
   indention=1.5em,
   labelfont={small,bf},
   up,
   justification=justified,
   singlelinecheck=true]{caption}

%% for sidewaysfigure and sidewaystable
\usepackage{rotating}

%% put recipe names, QC parameters, etc., in \lstinline{}
\usepackage{listings}
\lstloadlanguages{csh, c}

\usepackage{textcomp}

\usepackage{subfloat}

%% Write recipe names like this: \REC{metis_do_stuff}
\makeatletter
\lstdefinestyle{RECstyle}{%
  basicstyle=\ttfamily\color{darkgreen}%
  \lst@ifdisplaystyle\scriptsize\fi}
\makeatother
\newcommand{\REC}[1]{\lstinline[style=RECstyle]!#1!}

%% Write QC parameters like this: \QC{QC_SOMETHING_OR_OTHER}
\makeatletter
\lstdefinestyle{QCstyle}{%
  basicstyle=\ttfamily\color{darkblue}%
  \lst@ifdisplaystyle\scriptsize\fi}
\makeatother
\newcommand{\QC}[1]{\lstinline[style=QCstyle]!#1!}

%% Write templates like this: \TPL{DARK_LM}
\makeatletter
\lstdefinestyle{TPLstyle}{%
  basicstyle=\ttfamily\color{darkred}%
  \lst@ifdisplaystyle\scriptsize\fi}
\makeatother
\newcommand{\TPL}[1]{\lstinline[style=TPLstyle]!#1!}

%% Write products like this: \PROD{SOME_THING}
\makeatletter
\lstdefinestyle{PRODstyle}{%
  basicstyle=\ttfamily\color{darkorange}%
  \lst@ifdisplaystyle\scriptsize\fi}
\makeatother
\newcommand{\PROD}[1]{\lstinline[style=PRODstyle]!#1!}

%% Write requirements like this: \REQ{METIS-xxxx}
\newcommand{\REQ}[1]{\textcolor{brown}{#1}}

%% external calib files
\makeatletter
\lstdefinestyle{EXTCALIBstyle}{%
  basicstyle=\ttfamily\color{cyan}%
  \lst@ifdisplaystyle\scriptsize\fi}
\makeatother
\newcommand{\EXTCALIB}[1]{\lstinline[style=EXTCALIBstyle]!#1!}

% static calib files
\makeatletter
\lstdefinestyle{STATCALIBstyle}{%
  basicstyle=\ttfamily\color{teal}%
  \lst@ifdisplaystyle\scriptsize\fi}
\makeatother
\newcommand{\STATCALIB}[1]{\lstinline[style=STATCALIBstyle]!#1!}

%% Write FITS keywords (and values) like this: \FITS{EXPTIME}
\newcommand{\FITS}[1]{\lstinline[]!#1!}

%% Write code snippets like this: \CODE{SOMETHING==ELSE}
%% (I know, this is a bit stupid)
\newcommand{\CODE}[1]{\lstinline[]!#1!}

%% Use environment recipedef to define a new recipe.
\newenvironment{recipedef}
{% begin code
  \begin{tcolorbox}[breakable,enhanced jigsaw]%
    \begin{longtable}{p{0.24\hsize}p{0.74\hsize}}}
      {% end code
    \end{longtable}%
  \end{tcolorbox}}

%%%%% This is a definition that has the flow diagram within the box:
%% \newcommand{\theimage}{}  % dummy definition
%% \newenvironment{recipedef}[1]
%% {% begin
%% code
%% \renewcommand{\theimage}{#1}
%% \begin{tcolorbox}%
%%   \begin{tabular}{p{0.24\hsize}p{0.74\hsize}}
%%   }%
%%     {% end code
%%   \end{tabular}%
%%   \begin{center}%
%%     \includegraphics[width=0.6\textwidth]{\theimage}%
%%   \end{center}%
%% \end{tcolorbox}}


%% Proper typography for micrometer
\newcommand{\micron}{\textmu\mathrm{m}}

%% Ion
\DeclareRobustCommand{\ion}[2]{%
  \textup{#1{\mdseries\textsc{#2}}}%
}

%%%% \la...kleiner als ungef"ahr
\def\la{\mathrel{\mathchoice {\vcenter{\offinterlineskip\halign{\hfil
          $\displaystyle##$\hfil\cr<\cr\sim\cr}}}
    {\vcenter{\offinterlineskip\halign{\hfil$\textstyle##$\hfil\cr
          <\cr\sim\cr}}}
    {\vcenter{\offinterlineskip\halign{\hfil$\scriptstyle##$\hfil\cr
          <\cr\sim\cr}}}
    {\vcenter{\offinterlineskip\halign{\hfil$\scriptscriptstyle##$\hfil\cr
          <\cr\sim\cr}}}}}

%%%% \ga...Gr"o"ser als ungef"ahr
\def\ga{\mathrel{\mathchoice {\vcenter{\offinterlineskip\halign{\hfil
          $\displaystyle##$\hfil\cr>\cr\sim\cr}}}
    {\vcenter{\offinterlineskip\halign{\hfil$\textstyle##$\hfil\cr
          >\cr\sim\cr}}}
    {\vcenter{\offinterlineskip\halign{\hfil$\scriptstyle##$\hfil\cr
          >\cr\sim\cr}}}
    {\vcenter{\offinterlineskip\halign{\hfil$\scriptscriptstyle##$\hfil\cr
          >\cr\sim\cr}}}}}

%%%% Degrees, arcmin, arcsec
\def\degr{\hbox{$^\circ$}}

\def\arcmin{\hbox{$^\prime$}}

\def\arcsec{\hbox{$^{\prime\prime}$}}

%% Bright red for things that still need to be looked at. None of
%% these should appear in the release version
\newcommand{\TODO}[1]{\textcolor{red}{\bfseries TODO: #1}}

\newcommand{\intentblankpage}{}

%% paragraph column, centred text
\newcolumntype{P}[1]{>{\centering\arraybackslash}p{#1}}

%=== SUBVERSION INFO ==========================================================
\svnid{$Id$}
\newcommand{\svnrevisiondate}{\svnday-\svnmonth-\svnyear}  % <<<<< svn date

%=== Header / Footer definition ===============================================
\newcommand{\headerformat}{%
  \lhead{\small\sffamily
    \begin{tabularx}{\textwidth}{Xlll}
      \ifshowSVN
      \shorttitle & \docnumber & SVN \ \svnrev & \svnrevisiondate\\[0.7ex]
      \else
      \shorttitle & \docnumber & \issuenumber & \issuedate\\[0.7ex]
      \fi
    \end{tabularx}}
  \rhead{}
  \cfoot{\small\bfseries\sffamily
    Page {\textbf{\thepage}}  of {\textbf{\pageref{LastPage}}} }
  \rfoot{\raisebox{-0.3\height}{\includegraphics[width=25mm]{metis_noText.pdf}} }
}

\fancypagestyle{plain}{\fancyhf{} \headerformat}
\headerformat
\definecolor{brn}{RGB}{99,36,35}
\definecolor{rd1}{RGB}{229,184,183}
\definecolor{rd2}{RGB}{242,219,219}

\renewcommand\headrule{%
  \begingroup
  \color{brn}
  \hrule height 0.5pt width\headwidth
  \endgroup
}

\renewcommand{\contentsname}{Table of Contents}

%% ---- Format of section titles --------------------------------
% Combining sffamily and scshape works only for some fonts
\titleformat{\section}{\LARGE\sffamily\scshape}{\thesection}{1em}{\textsc{}}
\titleformat{\subsection}{\Large\sffamily}{\thesubsection}{1em}{}
\titleformat{\subsubsection}{\large\sffamily}{\thesubsubsection}{1em}{}

\setlength\heavyrulewidth{1.5pt}


% === PDF Definition ==========================================================
\hypersetup{%
  setpagesize,
  bookmarksnumbered=true,
  pdfborder= 0 0 0,
  pdftitle={\shorttitle},
  pdfauthor={\authorname}}
\pdfinfo{ /CreationDate (D:\svnpdfdate)}


% === biblatex Definitions ====================================================
\bibstyle{apj}
\addbibresource{references.bib}
\newbibmacro{string+doiurlisbn}[1]{\iffieldundef{url}{#1}{\href{\thefield{url}}{#1}}}
\DeclareFieldFormat{title}{\usebibmacro{string+doiurlisbn}{\mkbibemph{#1}}}


% === glossaries definitions  =================================================
% \input{../common/metis_glossary.tex}     % the common glossary file
\renewcommand{\glossarysection}[2][]{}   % Avoid the heading 'Glossary'
\setlength{\LTleft}{0pt}                 % make all longtables aligned left
\setlength{\glsdescwidth}{0.8\textwidth}


%%%%%%%%%%%%%%%%%%%%%%%%%%%%%%%%%%%%%%%%%%%%%%%%%%%%%%%%%%%%%%%%%%%%%%%%%%%%%
%
% Start of document
%
%%%%%%%%%%%%%%%%%%%%%%%%%%%%%%%%%%%%%%%%%%%%%%%%%%%%%%%%%%%%%%%%%%%%%%%%%%%%%
%
\begin{document}

% listings style
\lstset{basicstyle=\ttfamily\lst@ifdisplaystyle\scriptsize\fi,
  columns=flexible,
  frame=single,
  backgroundcolor=\color{listingbg},
  captionpos=b,
  showspaces=false}


%%%%% titlepage
\thispagestyle{empty}

% Logo
\vspace*{-1em}
\includegraphics[width=6.69cm]{metis_logo.pdf}

\vspace*{\fill}

% Document title
{\color{brn}\rule[1.9ex]{\textwidth}{1.5pt}}
\scalebox{1.44}{\Huge\textsf{\doctitle}}\\
{\color{brn}\rule{\textwidth}{1.5pt}}\\ [0.5ex]

% Document info
{\Large\textsf{\docnumber}}  \\ [1ex]
{\Large\textsf{Issue \issuenumber}} \\ [1ex]
{\Large\textsf{\issuedate}}  \\ [1ex]
{\Large\textsf{Work package: \workpackage}}  \\[1ex]

\vspace*{\fill}

% Signature table
\begin{center}
  \renewcommand{\arraystretch}{0.75}
%  \begin{tabularx}{\textwidth}{XXXX}
  \begin{tabular}{p{0.2\textwidth}p{0.4\textwidth}p{0.2\textwidth}p{0.2\textwidth}}
    \arrayrulecolor{brn}
    \toprule
    & \multicolumn{3}{l}{\scriptsize\textsf{Signature and Approval}} \\
    \midrule
    & {\scriptsize\textsf Name}
    & {\scriptsize\textsf Date}
    & {\scriptsize\textsf Signature} \\
      \cline{2-4} \\[-1mm]
    \textsf{Prepared} & \parbox[c]{\hsize}{\raggedright \authorname} & \authorsigndate & \\
    \midrule
    \textsf{Reviewed} & \parbox[c]{\hsize}{\raggedright \reviewername} & \reviewersigndate & \\
    \midrule
    \textsf{Approved} & \parbox[c]{\hsize}{\raggedright \approvername} & \approvalsigndate & \\
    \midrule
    \textsf{Released} & \parbox[c]{\hsize}{\raggedright \releasername} & \releasesigndate & \\
    \bottomrule
     %    \end{tabularx}
  \end{tabular}
\end{center}

%%% Local Variables:
%%% TeX-master: "METIS_DRLD"
%%% End:


% === MAIN TEXT================================================================

\clearpage
\pagestyle{fancy}

% === Revision History =========================================================
\section*{Revision History}

\renewcommand{\arraystretch}{1.2}
\arrayrulecolor{brn}
\begin{tabularx}{\textwidth}{|l|l|l|X|}
  \hline
  \rowcolor{rd1}
  \textbf{Issue}     & \textbf{Date} & \textbf{Owner} & \textbf{Changes}                           \\
  \hline
%  $10^{-6}$	     & 28-Feb-2017   &                & initial draft created                      \\
%  $1.01\cdot10^{-6}$  & 5-Apr-2017    &                & added HDRL-parameters                      \\
%                     &               &                & filled in recipes for N-band imaging       \\
%  $2\cdot10^{-6}$     & 9-May-2017    & none           & converted to official template format      \\
%  $3\cdot10^{-6}$     & 01-Jun-2017   & Kausch         & added Long-slit spectroscopy               \\
%  $4\cdot10^{-6}$     & 25-Aug-2017   & K\"ohler       & added IFU spectroscopy                     \\
%  $5\cdot10^{-6}$     & 5-Dec-2017    & K\"ohler       & added draft of instrument data description \\
  0.1                & 2020-05-27    & Czoske         & initial draft                             \\
  \hline
\end{tabularx}


%=== Indexes ==================================================================
\newpage
\tableofcontents
\clearpage
\listoffigures
\clearpage
\listoftables

\clearpage
\phantom{a}
\vfill
\begin{center}
  This page intentionally left blank
\end{center}
\vfill
\clearpage

%%%%%%%%%%%%%%%%%%%%%%%%%%%%%%%%%%%%%%%%%%%%%%%%%%%%%%%%%%%%%%%%%%%%%%%%%%%%%


% INTRODUCTION
\section{Introduction}
\label{sec:intro}

\subsection{Scope}

This document describes the design of the data reduction software for
METIS, the Mid-infrared ELT Imager and Spectrograph. It is based on
the Data Reduction Library Specifications document \cite{DRLS} presented
for PDR.

\subsection{Applicable documents}

\begin{refcontext}[labelprefix=AD]
  \printbibliography[keyword=applicable, heading=none]
\end{refcontext}


\subsection{Reference documents}

\begin{refcontext}[labelprefix=RD]
  \printbibliography[keyword=reference, heading=none]
\end{refcontext}


\clearpage

\subsection{Acronyms}
\label{ssec:acronyms}
%%% Acronym list

This document employs several abbreviations and acronyms to refer concisely to an item, after it has been introduced.
The following list is aimed to help the reader in recalling the extended meaning of each short expression:

\begin{acronym}[AAAAAAA]
   \acro{ADC}{Atmospheric Dispersion Corrector}
   \acro{ADI}{Angular Differential Imaging}
   \acro{ADU}{Analogue-Digital-Unit}
   % \acro{AGB}{Asymptotic Giant Branch}
   \acro{AIT}{Assembly Integration and Test}
%    \acro{AIV}{Assembly Integration and Verification}
   \acro{AO}{Adaptive Optics}
%   \acro{APP}{Apodizing Phase Plate}
%   \acro{CFO}{Common Fore Optics}
%   \acro{CLC}{Classical Lyot Coronagraph}
   \acro{CRIRES}{CRyogenic high-resolution InfraRed Echelle Spectrograph}
   \acro{CPL}{Common Pipeline Library}
%   \acro{CVC}{Classical Vortex Coronagraph}
%   \acro{DFS}{Data Flow System}
%   \acro{DIT}{Detector Integration Time}
   \acro{DRL}{Data Reduction Library}
%   \acro{DRLD}{Data Reduction Library Design}
   \acro{DRS}{Data Reduction Software}
   \acro{EDPS}{ESO Data Processing System}
   \acro{ELFN}{Excess Low Frequency Noise}
   \acro{ELT}{Extremely Large Telescope}
   \acro{ESO}{European Southern Observatory}
   \acro{FDR}{Final Design Review}
   \acro{FITS}{Flexible Image Transport System}
   \acro{FWHM}{Full Width Half Maximum}
%   \acro{FP}{Focal Plane}
   \acro{HCI}{High-Contrast Imaging}
   \acro{HDRL}{High-Level Data-Reduction Library}
   \acro{HITRAN}{HIgh-resolution TRANsmission molecular absorption database}
   \acro{HST}{Hubble-Space-Telescope}
%   \acro{HDU}{Header Data Unit}
   \acro{ICS}{Instrument Control System}
%   \acro{IFS}{Integral Field Spectrograph}
   \acro{IAU}{International Astronomical Union}
   \acro{IFU}{Integral Field Unit}
   \acro{IMG}{Imaging Mode}
%   \acro{IR}{Infrared}
   \acro{IRTF}{NASA Infrared Telescope Facility}
   \acro{JWST}{James-Webb-Space-Telescope}
   \acro{KMOS}{K-band Multi Object Spectrograph}
   \acro{LHATPRO}{Humidity And Temperature PROfiler}
   \acro{LSF}{Line Spread Function}
   \acro{LMS}{L/M band Spectrograph Subsystem}
   \acro{LSS}{Long-Slit Spectroscopy}
%   \acro{MAIT}{Manufacture, Assembly, Integration and Test}
   \acro{METIS}{Mid-infrared E-ELT Imager and Spectrograph}
   \acro{MIR}{Mid-Infrared}
%   \acro{NIR}{Near-Infrared}
   \acro{OB}{Observation Block}
   \acro{PAE}{Preliminary Acceptance Europe}
   \acro{PDR}{Preliminary Design Review}
%   \acro{PP}{Pupil Plane}
   \acro{PSF}{Point Spread Function}
   \acro{PWV}{Precipitable Water Vapour}
   \acro{QC}{Quality Control}
   \acro{QCL}{Quantum Cascade Laser}
%   \acro{RAVC}{Ring Apodized Vortex Coronagraph}
   \acro{RMS}{Root Mean Square}
   \acro{RSRF}{Relative Spectral Response Function}
   \acro{SNR}{Signal-to-Noise Ratio}
      \acro{SCAO}{Single Conjugate Adaptive Optics}
%   \acro{SED}{Spectral Energy Distribution}
   \acro{SDP}{Science Data Product}
   \acro{SOF}{Set-of-Files}
   \acro{STD}{Standard star}
%   \acro{TBC}{To Be Confirmed}
%   \acro{TBD}{To Be Determined}
%   \acro{VLT}{Very Large Telescope}
    \acro{TSS}{Telluric Standard Star}
%   \acro{WCS}{World Coordinate System}
    \acro{WFS}{Wave Front Sensor}
  \acro{VISIR}{VLT Imager and Spectrometer for mid Infrared}
  \acro{WCU}{Warm Calibration Unit}
\end{acronym}

%%%%%%%%%%%%%%%%%%%%%%%%%%%%%%%%%%%%%%%%%%%%%%%%%%%%%%%

%%% Local Variables:
%%% TeX-master: "METIS_DRLD"
%%% End:


\clearpage

% \subsection{Specifications and requirements}
% \label{Subsec:polarion}
% \input{01-Requirements}


\end{document}

%%
%% THE END
%%
%%%%%%%%%%%%%%%%%%%%%%%%%%%%%%%%%%%%%%%%%%%%%%%%%%%%%%%%%%%%%%%%%%%%%%%%%%%%%
